\section{Introduction}
\label{sec:introduction}

Mean dimension measures the dimension needed per orbit coordinate.
Introduced by Gromov and developed by Lindenstrauss and Weiss, it remains
informative for compact dynamical systems whose ordinary covering
dimension is infinite \cite{Gromov1999,LindenstraussWeiss2000}.
Toms introduced the radius of comparison to measure the tracial gap
needed for comparison of positive elements of a C*-algebra
\cite[Definition~6.1]{Toms2006}. The Phillips--Toms conjecture relates
these dynamical and operator-algebraic quantities. We prove its
integer-action case, with the factor one half suggested by the relation
between complex vector-bundle rank and real dimension.

\subsection{The statement and conventions}

For a finite open cover \(\mathcal V\) of a compact space, let
\(D(\mathcal V)\) be the least order of a finite open refinement, where
the order is one less than the largest number of members with a nonempty
common intersection. For a homeomorphism \(h:X\to X\), put
\[
\mathcal U^{[a,b)}=\bigvee_{j=a}^{b-1}h^{-j}\mathcal U,
\qquad
\mdim(X,h)=\sup_{\mathcal U}
\lim_{N\to\infty}\frac{D(\mathcal U^{[0,N)})}{N}.
\]
The supremum ranges over finite open covers. The limit for each cover
exists by subadditivity and is finite; the supremum may be infinite.

We use the tracial formulation of the radius of comparison for nuclear
C*-algebras \cite[Definition~2.10]{Niu2024}; in this setting quasitraces
are traces. Let \(A\) be a unital nuclear C*-algebra with tracial states. Write
\(M_\infty(A)=\bigcup_{q\ge1}M_q(A)\), using the upper-left corner
inclusions. For positive elements \(a,b\in M_\infty(A)\),
\(a\precsim b\) denotes Cuntz subequivalence: there are matrices
\(x_j\), after passage to compatible finite corners, such that
\(x_j^*bx_j\to a\) in norm. For a normalized tracial state \(\tau\), set
\[
d_\tau(a)=\lim_{m\to\infty}(\tau\otimes\operatorname{Tr})(a^{1/m}),
\]
where \(\operatorname{Tr}\) is the unnormalized matrix trace.
We say that a finite \(r\ge0\) is a \emph{comparison bound} for \(A\) if
\begin{equation}
\label{eq:comparison-bound}
\bigl(d_\tau(a)+r<d_\tau(b)\ \text{for every normalized trace }\tau\bigr)
\quad\Longrightarrow\quad a\precsim b
\end{equation}
for all such positive elements. The radius of comparison \(\rc(A)\)
is the infimum of the comparison bounds, with \(\inf\varnothing=\infty\).

\begin{theorem}
\label{thm:main}
Let \(h\) be a minimal homeomorphism of an infinite compact metrizable
space \(X\). Then
\begin{equation}
\label{eq:main}
\rc\bigl(C(X)\rtimes_h\Z\bigr)=\frac12\mdim(X,h)
\qquad\text{in }[0,\infty].
\end{equation}
\end{theorem}

Infinitude and minimality imply that the action is free: a periodic orbit
would be a finite closed invariant subset equal to \(X\).
The crossed product is nuclear \cite[Section~3]{ElliottNiu2017} and
has tracial states, obtained from invariant probability measures and
the canonical conditional expectation.
Thus the conventions above apply. In particular, \Cref{thm:main}
asserts that infinite mean dimension forces infinite radius of
comparison, as well as identifying the exact constant at finite
positive mean dimension.

\subsection{Prior results and the remaining lower bound}

The two invariants arose in different settings. Toms's radius of
comparison records the dimension-to-rank obstruction to comparing
positive elements, extending the role of vector-bundle obstructions
in C*-algebra classification \cite[Section~6]{Toms2006}.
Giol--Kerr brought this obstruction into dynamics: they constructed
minimal systems of positive mean dimension whose crossed products have
arbitrarily large radii of comparison
\cite[Section~2 and Theorem~2.2]{GiolKerr2010}.
The Phillips--Toms conjecture predicts an exact quantitative relation.
Its integer-action formulation is recorded in Elliott--Niu
\cite[Introduction]{ElliottNiu2017} and as Problem~XXXVI in
Schafhauser--Tikuisis--White \cite{SchafhauserTikuisisWhite2026}.
Niu states the broader formulation for free minimal actions of discrete
amenable groups \cite[Introduction]{Niu2024}. Our theorem concerns
minimal homeomorphisms of infinite compact metrizable spaces.

The upper bounds developed from a coarse dimension estimate to the
conjectured constant. Phillips proved
$\rc(C(X)\rtimes_h\Z)\le1+36\mdim(X,h)$
\cite[Corollary~4.8]{Phillips2016}, while Elliott--Niu settled the
zero-mean-dimension case \cite[Theorem~4.7]{ElliottNiu2017}.
Niu then proved the optimal upper bound
\begin{equation}
\label{eq:known-upper-bound}
\rc\bigl(C(X)\rtimes_h\Z\bigr)\le\tfrac12\mdim(X,h)
\end{equation}
for free minimal integer-lattice actions, including the systems
considered here \cite[Theorem~5.6]{Niu2024}. Thus equality at zero
mean dimension is already known. Niu's open tower theorem, with an
arbitrarily small uniform orbit remainder, is a second input to our
proof \cite[Theorem~4.2]{Niu2024}.

Hirshberg--Phillips obtained lower bounds for actions of countable
amenable groups using mean cohomological independence, which detects
nonzero cup products of translated cohomology classes
\cite[Theorems~3.3 and~4.5 in the author preprint]{HirshbergPhillips2022}.
For the minimal subshifts constructed by Dou \cite[Section~4]{Dou2017} over positive
even-dimensional polyhedra with nonzero top rational cohomology, their
bound is strictly greater than one half of the mean dimension minus one
\cite[Corollary~4.6 in the author preprint]{HirshbergPhillips2022}.
Their arguments turn finite-support comparison witnesses into bundle
embeddings and obstruct small target rank with inverse Chern classes.

Covering dimension itself can also yield comparison lower bounds:
Phillips proved $\rc(C(Y))\ge(\dim Y-7)/2$ for every compact Hausdorff
space $Y$, using extension of sphere-valued maps
\cite[Corollary~3.4]{Phillips2023}.
This commutative result does not control comparison in the larger
crossed product. The remaining dynamical problem is to recover every
orbit-cover rate from crossed-product comparison. Our argument retains
the connection between comparison and bundles, but converts the bundle
information directly into a reduction in covering order. It therefore
applies to every finite open cover, without a cohomological independence
hypothesis or a finite-dimensionality hypothesis on $X$.

\subsection{Zero mean dimension and regularity}

For \(E\subset X\), define its orbit capacity by
\[
\operatorname{ocap}_h(E)=\lim_{N\to\infty}\frac1N\sup_{x\in X}
   \sum_{j=0}^{N-1}\mathbf 1_E(h^jx).
\]
The limit exists by subadditivity. The system \((X,h)\) has the
\emph{small boundary property} if every \(x\in X\) and every open
neighborhood \(U\) of \(x\) admit an open neighborhood \(V\) of \(x\)
with \(V\subset U\) and \(\operatorname{ocap}_h(\partial V)=0\).
A C*-algebra \(A\) is \emph{Jiang--Su stable} if
\(A\cong A\otimes\mathcal Z\), where \(\mathcal Z\) is the Jiang--Su
algebra and the tensor product is spatial. Nuclear dimension
\(\dim_{\mathrm{nuc}}(A)\) is defined through finite-dimensional
completely positive approximations, with a bound \(n\) permitting
\(n+1\) contractive order-zero return pieces. An order-zero map
preserves orthogonality of positive elements; the nuclear dimension is
\(\infty\) when there is no finite bound.

The zero endpoint has a further consequence: it characterizes the
regularity of the crossed product. The dynamical and operator-algebraic
conditions in this characterization use different established inputs.

\begin{corollary}[Zero mean dimension and crossed-product regularity]
\label{cor:regularity}
Let \(h\) be a minimal homeomorphism of an infinite compact metrizable
space \(X\), and put \(A=C(X)\rtimes_h\Z\). The following conditions are
equivalent:
\begin{enumerate}
\item \(\mdim(X,h)=0\);
\item \((X,h)\) has the small boundary property;
\item \(A\cong A\otimes\mathcal Z\);
\item \(\dim_{\mathrm{nuc}}(A)<\infty\);
\item \(\dim_{\mathrm{nuc}}(A)\le1\).
\end{enumerate}
In particular, if \(\mdim(X,h)\in(0,\infty]\), then
\(A\not\cong A\otimes\mathcal Z\) and
\(\dim_{\mathrm{nuc}}(A)=\infty\).
\end{corollary}

The equivalence of zero mean dimension and the small boundary property
is due to Lindenstrauss--Weiss and Lindenstrauss
\cite[Theorem~5.4]{LindenstraussWeiss2000}
\cite[Theorem~6.2]{Lindenstrauss1999}; the latter theorem applies because
our system is itself an infinite minimal factor.
Elliott--Niu proved that zero mean dimension implies Jiang--Su stability
\cite[Theorem~4.7]{ElliottNiu2017}. Conversely, Jiang--Su stability gives
comparison at zero by R\o rdam's theorem
\cite[Corollary~4.6]{Rordam2004}, and \Cref{thm:main} then forces zero
mean dimension. The remaining equivalences use the published
simple-algebra nuclear-regularity theorems
\cite[Theorems~A and~B]{CastillejosEvingtonTikuisisWhiteWinter2021}.
We verify their hypotheses and the matrix-trace conventions in
\Cref{sec:regularity-proof}.

\subsection{The proof: comparison, compression, and colors}

Fix a finite comparison bound $r$ and a finite open cover $\mathcal U$.
The main task is to prove that the cover rate
\[
 d_{\mathcal U}=\lim_{k\to\infty}
            \frac{D(\mathcal U^{[0,k)})}{k}
\]
is at most $2r$. Assume instead that $d_{\mathcal U}>2r$.
A subordinate nerve represents one long orbit block by a point in a
finite simplicial complex of dimension $m=D(\mathcal U^{[0,k)})$.
After barycentric subdivision, its vertices can be colored with $m+1$
colors, distinct on every simplex. The barycentric weight of a color
is the sum of the weights of vertices of that color.

For a large positive integer $p$, group colors into packets of $2p+2$.
Their proportions give points in a $(2p+1)$-dimensional simplex, which
we identify with a cube. Since proportions are undefined at zero packet
mass, we regularize them to obtain continuous cube-valued maps. We call
each cube factor a \emph{slot}. To replace slot maps while retaining
their cover information, we must preserve every originally absent color. The compression theorem
from the companion paper
\emph{Filtered products and boundary-preserving compression in complex
cobordism} \cite{OpenAICompression2026} supplies precisely the needed
control: after passing to arbitrarily large product powers, it puts a
positive fraction of the slots on their boundaries and fixes the entire
vector in every original boundary slot. Its full statement and bundle
hypothesis appear in \Cref{sec:compression-interface}.

Here is how comparison supplies that hypothesis. In each slot fix a
small framed torus $(\mathbb T^2)^p$ and a rank-$p$ bundle formed from
one degree-one line bundle on each two-torus factor. The compression
hypothesis tests the slot maps after translations by
parameters supported in only a small fraction of the slots. The resulting
torus loci carry a sum of these rank-$p$ bundles.
In \Cref{sec:dynamical-comparison}, an open tower allocates orbit blocks
with continuous alignment labels on most starting points. A single
finite-propagation Cuntz witness embeds the sum of the aligned bundle
fibers. Two fixed linear maps handle the varying coordinate supports,
including sparse parameters and the few blocks without full alignment.
This produces one continuous embedding over the entire parameter locus,
with target rank less than twice the source rank
(\Cref{prop:comparison-bundle}). The continuity through changes of sparse
support is needed by the compression theorem.

It remains to interpret its output as an actual smaller cover.
\Cref{sec:cover-reconstruction} proves a packet reconstruction lemma.
It discards packets of sufficiently small total weight; on retained
packets, exact boundary preservation prevents any originally absent
color from appearing. Together these operations preserve every original
zero vertex coordinate, while each boundary output leaves at least one
color absent from its packet. \Cref{sec:lower-bound} applies this
reconstruction on aligned orbit blocks. The two edge blocks retain the beginning and
end of each interval; a join with an unchanged description handles
points where alignment disappears. A count of the cells containing the
resulting image places it in one finite complex of smaller dimension.
Its open-star cover then refines the entire orbit cover, contradicting
the rate $d_{\mathcal U}$.

The conversion from $p$ complex lines to $2p+2$ colors has ratio
$p/(2p+2)$, which tends to one half. We choose all tower, sparse-parameter,
rounding, and compression losses within the initial strict margin
$d_{\mathcal U}>2r$. In particular, one orbit-window length is fixed
before invoking the arbitrarily large product powers. The final passage
takes the supremum over covers and the infimum over comparison bounds;
when mean dimension is infinite, the individual-cover bound excludes
every finite comparison bound. Niu's upper bound then completes
\Cref{thm:main}.
